\documentclass[12pt,a4paper]{article}

\usepackage{amsmath, amssymb}
\usepackage{geometry}
\geometry{margin=2.5cm}

\title{Modelarea Matematică a Sistemelor Materiale \\ \large Capitolul 9: Calcul Tensorial}
\author{}
\date{}

\begin{document}

\maketitle

\section{Calcul tensorial}

Calculul tensorial este folosit în mecanica teoretică pentru a descrie relațiile dintre forțe,
mişcări, tensiuni și deformații într-un mod independent de sistemul de coordonate.

\subsection{9.1. Definiția tensorului}

Un tensor de ordinul doi este o aplicație liniară $T : V \to V$ care transformă un vector $u$
într-un vector $v$:


\[
T(\alpha u + \beta v) = \alpha Tu + \beta Tv.
\]



Doi tensori sunt egali dacă:


\[
Tu = Su, \quad \forall u \in V.
\]



\subsection{9.2. Componentele unui tensor}

Fie $\{i_1,i_2,i_3\}$ o bază ortonormată. Tensorul $T$ aplicat pe vectorii bazei produce:


\[
Ti_j = T_{1j} i_1 + T_{2j} i_2 + T_{3j} i_3.
\]



Matricea asociată tensorului este:


\[
[T] =
\begin{pmatrix}
T_{11} & T_{12} & T_{13} \\
T_{21} & T_{22} & T_{23} \\
T_{31} & T_{32} & T_{33}
\end{pmatrix}.
\]



Componentele se obțin prin:


\[
T_{ij} = i_i \cdot (T i_j).
\]



\subsection{9.3. Produsul tensorial}

Pentru $u,v \in V$, produsul tensorial este aplicația:


\[
(u \otimes v) w = u (v \cdot w).
\]



Tensorii $i_i \otimes i_j$ formează o bază a spațiului $\mathrm{Lin}$.

\subsection{9.4. Operații cu tensori}

\paragraph{Suma}


\[
(T+S)_{ij} = T_{ij} + S_{ij}.
\]



\paragraph{Produsul de compoziție}


\[
(TS)_{ij} = T_{ik} S_{kj}.
\]



\paragraph{Transpusul}


\[
(T^T)_{ij} = T_{ji}.
\]



\paragraph{Urma}


\[
\mathrm{tr}(T) = T_{11} + T_{22} + T_{33}.
\]



\paragraph{Determinantul}


\[
\det T = \det [T].
\]



\paragraph{Produsul scalar}


\[
T \cdot S = T_{ij} S_{ij}.
\]



\paragraph{Norma}


\[
\|T\| = \sqrt{T \cdot T}.
\]



\subsection{9.5. Schimbarea sistemului de coordonate}

Dacă $Q$ este un tensor ortogonal ($Q^T Q = I$), atunci:



\[
[a]' = Q^T [a],
\]




\[
[T]' = Q^T [T] Q.
\]



\subsection{9.6. Tensori simetrici și antisimetrici}

\paragraph{Tensor simetric}


\[
T^T = T.
\]



\paragraph{Tensor antisimetric}


\[
T^T = -T.
\]



Orice tensor se descompune în:


\[
T = T_S + T_A,
\]




\[
T_S = \frac{1}{2}(T + T^T), \quad
T_A = \frac{1}{2}(T - T^T).
\]



\subsection{9.7. Vectorul coaxial al unui tensor antisimetric}

Pentru un tensor antisimetric $\Omega$ există un vector $\omega$ astfel încât:


\[
\Omega v = \omega \times v.
\]



Componentele sunt:


\[
\omega_p = -\frac{1}{2} \varepsilon_{pki} \Omega_{ki}.
\]



\subsection{9.8. Exemple fundamentale}

\paragraph{Tensorul rotației}
Pentru rotație cu unghi $\theta$ în jurul vectorului unitar $n$:


\[
R = (1-\cos\theta)(n \otimes n) + \cos\theta\, I + \sin\theta\, \Omega_n.
\]



\paragraph{Tensorul reflexiei}
Pentru un plan cu normală unitară $n$:


\[
T = I - 2 n \otimes n.
\]



\end{document}